\documentclass[11pt]{article}
\usepackage[margin=2.5cm]{geometry}
\usepackage[T1]{fontenc}
\usepackage[utf8]{inputenc}
\setlength{\parskip}{0.7em}
\setlength{\parindent}{0pt}
\pagestyle{empty}
\begin{document}
To the Editors, The European Physical Journal B

Dear Editors,

We submit ``Non-Monotonic Correlation Fluctuations in Heisenberg Models on Small Quantum Graphs'' for consideration as a regular article.

\textbf{Summary and main results.} We study the isotropic Heisenberg model on small graphs in which one edge, the contradiction edge, carries a coupling \texttt{s} that is scanned while every other edge stays at unity. The observable is the spatial standard deviation of nearest-neighbour correlations across all edges, which is relational rather than site-centred. On a level-2 Sierpinski gasket with N = 15 the observable falls into a sharp valley at the uniform point s = 1, reproducible across five Lanczos seeds; a degeneracy control shows the valley survives at non-zero gap. The result we consider the most useful is not the valley itself but its floor. We show that the valley floor is fixed by the edge automorphism orbit structure of the uniform graph: for any symmetry-invariant ground state the within-orbit variance vanishes identically, so the observable at s = 1 reduces to its between-orbit component and is exactly zero for edge-transitive graphs. The argument requires no thermodynamic limit, and it has been checked numerically on thirty distinct graphs, with random-block and symmetry-breaking controls that separate the orbit structure from the numerical distribution of correlations. The paper also reports a clean negative result for the PXP model, and it withdraws two earlier claims of our own: an apparent jump in the PXP scan, traced to a non-Hermitian Hamiltonian in an early script, and an earlier truncation figure that did not reproduce.

\textbf{Potential impact.} The paper makes a concrete and testable statement about when a fine-grained relational diagnostic can and cannot see structure that site-averaged observables miss. The symmetry constraint is a sufficient condition, stated as such, and we are explicit that it is not necessary. We expect the orbit-resolved argument to be reusable on other graph geometries, and we have deliberately kept the scope of the phenomenological part narrow: the manuscript states in its introduction that it is a phenomenology report and claims neither a new physical mechanism nor universality.

\textbf{Why EPJ B.} The work sits in condensed matter and statistical physics: it is an exact diagonalization and tensor-renormalization-group study of correlation structure in quantum spin models on graphs, and its central result is a symmetry statement about ground-state correlations. EPJ B publishes precisely this combination of numerical condensed-matter work and structural argument, and its readership in quantum magnetism and complex networks is the audience best placed to test the orbit-resolved claim on geometries we did not reach. The negative PXP result and the two withdrawn claims are also, we think, appropriate to a journal that publishes complete accounts rather than only positive findings.

\textbf{Data, code, and reproducibility.} All data and all scripts are public, with the recorded output of each run, in the two repositories named in the data availability statement. Every number attributed to the independent verification can be re-derived by a referee without contacting us.

\textbf{Artificial intelligence involvement.} Part of the verification was carried out with an AI research system, and this is documented in the Method section as the journal's policy requires. The system is not an author. The named authors are accountable for the whole of the text and for every number in it.

\textbf{Prior work on this lattice.} The manuscript cites the four earlier studies of the quantum Heisenberg antiferromagnet on the Sierpinski gasket, including the 2004 study in this journal. None of them examines the response to a single non-uniform edge, which is the question here.

\textbf{Suggested reviewers and exclusions.} We suggest two referees, both entered in the submission form: Haiyuan Zou, author of the most recent study of this model, Chin. Phys. Lett. 40, 057501 (2023); and Piotr Tomczak, co-author of the three earlier studies of the quantum Heisenberg antiferromagnet on the Sierpinski gasket, including Eur. Phys. J. B 38, 49 (2004) in this journal. We request no exclusions.

The manuscript has not been published elsewhere and is not under consideration by any other journal.

Yours sincerely,

Rastislav Drahos, corresponding author, on behalf of Guanghao Li and Marat Sultanov
\end{document}
